\addcontentsline{toc}{subsubsection}{Solution 57. The Hidden Interval}

\begin{dlsummary}[title={\bfseries Solution 57. The Hidden Interval},breakable]

\begin{proof}
The solutions are precisely the functions that \textbf{move each real number to the nearest point of a nonempty closed interval}, allowing unbounded intervals and single points.

Explicitly:

| Fixed-point interval | Function |
|---|---|
| \([a,b]\), where \(a\le b\) | \(\displaystyle f(x)=\begin{cases}a,&x<a,\\x,&a\le x\le b,\\b,&x>b;\end{cases}\) |
| \([a,\infty)\) | \(f(x)=\max\{x,a\}\) |
| \((-\infty,b]\) | \(f(x)=\min\{x,b\}\) |
| \(\mathbb R\) | \(f(x)=x\) |

The case \(a=b\) includes every constant function.

\paragraph{1. Extract order and continuity.}

For \(x>y\), put \(d=f(x)-f(y)\). Condition (1) says
\[
d^2\le(x-y)d.
\]
Equivalently,
\[
d\bigl(d-(x-y)\bigr)\le0,
\]
so
\[
\boxed{0\le f(x)-f(y)\le x-y.} \tag{3}
\]
Thus \(f\) is nondecreasing and
\[
|f(x)-f(y)|\le|x-y|.
\]
In particular, \(f\) is continuous.

\paragraph{2. Collapse the threefold iteration.}

Fix \(x\), and write
\[
u=f(x),\qquad v=f(u).
\]
By (2),
\[
f(v)=u.
\]

If \(u<v\), monotonicity gives
\[
v=f(u)\le f(v)=u,
\]
a contradiction. If \(u>v\), the same argument gives the opposite contradiction. Hence \(u=v\).

Therefore
\[
\boxed{f(f(x))=f(x).} \tag{4}
\]

\paragraph{3. Identify the fixed-point set.}

Let
\[
C=\{t\in\mathbb R:f(t)=t\}.
\]
Equation (4) shows that every function value belongs to \(C\), so \(C\) is nonempty and equals the image of \(f\). Continuity makes \(C\) closed.

Suppose \(a,b\in C\) and \(a<t<b\). Applying (3) between \(a,t\),
\[
f(t)-a\le t-a,
\]
so \(f(t)\le t\). Applying it between \(t,b\),
\[
b-f(t)\le b-t,
\]
so \(f(t)\ge t\).

Thus \(f(t)=t\). Consequently, \(C\) is an interval.

\paragraph{4. Determine the function outside the interval.}

If \(C\) has a finite lower endpoint \(a\), then \(a\in C\). For \(x<a\), monotonicity gives \(f(x)\le f(a)=a\), while \(f(x)\in C\) gives \(f(x)\ge a\). Hence
\[
f(x)=a.
\]

Similarly, above a finite upper endpoint \(b\), we have \(f(x)=b\). On \(C\), the function is the identity. This gives exactly the listed formulas.

\paragraph{5. Verify.}

Each listed function is nondecreasing and satisfies
\[
0\le f(x)-f(y)\le x-y\qquad(x>y),
\]
which implies (1). Its values lie in an interval that it fixes pointwise, so \(f(f(x))=f(x)\), proving (2).
\end{proof}

\paragraph{Takeaways.}
\begin{itemize}
\item **Factor a quadratic inequality:** It can simultaneously reveal monotonicity and a distance bound.
\item **Order rules out nontrivial two-cycles:** This turns threefold iteration into idempotence.
\item **Study the fixed-point set:** Its geometry can determine the entire function.
\item **Expect piecewise answers:** A functional inequality need not force a polynomial or power function.
\end{itemize}

\end{dlsummary}
